% Generated draft: requires mathematical and source review.
\documentclass[11pt]{article}
\usepackage[margin=25mm]{geometry}
\usepackage{fontspec}
\setmainfont{Latin Modern Roman}
\usepackage{amsmath,amssymb,mathtools,unicode-math}
\setmathfont{Latin Modern Math}
\usepackage{enumitem,needspace,xurl,hyperref}
\setlength{\parindent}{0pt}
\setlength{\parskip}{5pt}
\setlength{\emergencystretch}{3em}
\newcommand{\N}{\mathbb{N}}
\newcommand{\Z}{\mathbb{Z}}
\newcommand{\Q}{\mathbb{Q}}
\newcommand{\R}{\mathbb{R}}
\newcommand{\C}{\mathbb{C}}
\newcommand{\F}{\mathbb{F}}
\newcommand{\A}{\mathbb{A}}
\newcommand{\PP}{\mathbb{P}}
\newcommand{\E}{\mathbb{E}}
\newcommand{\eps}{\varepsilon}
\newcommand{\dd}{\,\mathrm{d}}
\begin{document}
\section*{UnsolvedMath ID: 1099}
\textbf{Source problem code:} catalogue entry
\subsection{Definitions and mathematical statement}
For $X\in\N$ write $[X]=\{1,\ldots,X\}$. For $n\ge1$, an integer polynomial $F\in\Z[x_1,\ldots,x_n]$, an integer $C$, and $B\subseteq[X]$, put
\[
 N_{F,C}(B)=\bigl|\{(x_1,\ldots,x_n)\in B^n:F(x_1,\ldots,x_n)=C\}\bigr|.
\]
These are ordered tuples; no distinctness restriction is imposed by the supplied statement. For $A_X\subseteq[X]$, let $\alpha_X=|A_X|/X$. The expected-count programme asks which hypotheses on $A_X$, on $F$, and on $C$ guarantee a relation of the form
\[
 N_{F,C}(A_X)=\alpha_X^nN_{F,C}([X])
                 +o\bigl(N_{F,C}([X])\bigr),\qquad X\longrightarrow\infty,
\]
in a fixed positive density regime, for example $\alpha_X\longrightarrow\alpha\in(0,1]$. When $N_{F,C}([X])$ is eventually positive, this is equivalently the ratio tending to $\alpha^n$. This display makes one natural meaning of ``roughly'' explicit; the source record does not specify the error scale or all required quantifiers, so it is not being asserted as an already precise universal conjecture. Eventual vanishing of the comparison count would require a separate convention. The record also does not specify the allowed dependence of $F,C$ on $X$, or the exact coefficient, nonsingularity, and local-solubility assumptions. Integer coefficients and integer $C$ describe the Diophantine setting of the supplied discussion and must be checked against the original problem if a precise narrower target is adopted.

The supplied triage identifies three concrete directions in Green's source: generic quadratic forms in seven variables; a useful least variable count for generic cubic forms; and the four-square equation under sharp pseudorandomness conditions on $A_X$. It does not reproduce the definition of ``generic,'' the full cubic quantifiers, the precise four-square equation, or the sharp pseudorandomness hypotheses. Those targets therefore cannot be uniquely formalized from this record alone.

As dated partial progress, the record reports that ordinary Fourier uniformity controls expected counts for linear equations with at least three variables; that Cook--Magyar provide high-variable local-to-global counting for polynomial systems with prime coordinates; that Green proves an asymptotic for mildly nondegenerate quadratic forms in eight prime variables; and that Green's maintained notes say the latter methods should yield an eight-variable generic dense-set result under non-correlation with progressions. The last statement is a proposed methodological extension in those notes, not a theorem asserted here. Browning--Prendiville give dense-set transference for translation-invariant diagonal quadrics in at least five variables within the squares. The record does not supply the precise hypotheses of these comparison theorems, and they are not universal assertions about every $F$ or $A_X$.

The registry's main label is open, while its literature triage checked 17 August 2026 classifies partial progress and reports that no general criterion or resolution of the three specific targets was located. It regards the stored question as an underspecified research programme rather than a binary proposition.
\subsection{Short English statement}
Seek hypotheses ensuring that a dense subset of $[X]$ contains the expected density power of all solutions to a polynomial equation. The supplied record reports nearby counting theorems but leaves the general hypotheses and the precise seven-variable quadratic, cubic, and four-square targets incomplete.
\subsection{Sources}
\begin{itemize}
\item Local UnsolvedMath registry, permanent ID $1099$, source code \texttt{catalogue entry}, ``Polynomial Solutions in Dense Sets,'' including the literature triage checked 17 August 2026.
\item Ben Green, \emph{100 Open Problems}, Problem 99, updates through December 2025 according to the record, \url{https://people.maths.ox.ac.uk/greenbj/papers/open-problems.pdf}.
\item Ben Green, \emph{Quadratic forms in 8 prime variables}, Geometric and Functional Analysis \textbf{35} (2025), 1587--1637, \url{https://doi.org/10.1007/s00039-025-00727-9}.
\item Brian Cook and Akos Magyar, \emph{Diophantine equations in the primes}, Inventiones Mathematicae \textbf{198} (2014), 701--737, \url{https://doi.org/10.1007/s00222-014-0508-1}.
\item Tim Browning and Sean Prendiville, \emph{A transference approach to a Roth-type theorem in the squares}, International Mathematics Research Notices 2017, 2219--2248, \url{https://doi.org/10.1093/imrn/rnw096}.
\end{itemize}
\end{document}
